\documentclass[11pt]{article}
\usepackage{amsmath,amssymb,amsthm}
\usepackage{hyperref}
\usepackage[margin=1in]{geometry}

\newcommand{\N}{\mathbb{N}}
\newcommand{\Z}{\mathbb{Z}}
\newcommand{\Q}{\mathbb{Q}}
\newcommand{\La}{\Lambda}
\newcommand{\GK}{\operatorname{GK}}
\newcommand{\CPP}{\operatorname{CPP}}
\newcommand{\ALCD}{\operatorname{ALCD}}
\newcommand{\IP}{\operatorname{IP}}
\newcommand{\qpoch}[2]{(#1;#2)_\infty}
\newcommand{\qpochn}[3]{(#1;#2)_{#3}}
\newcommand{\qbin}[2]{\genfrac{[}{]}{0pt}{}{#1}{#2}}
\DeclareMathOperator{\tr}{tr}

\theoremstyle{plain}
\newtheorem{theorem}{Theorem}
\newtheorem{conjecture}[theorem]{Conjecture}
\newtheorem{proposition}[theorem]{Proposition}
\newtheorem{lemma}[theorem]{Lemma}
\newtheorem{corollary}[theorem]{Corollary}
\theoremstyle{definition}
\newtheorem{definition}[theorem]{Definition}
\newtheorem{example}[theorem]{Example}
\newtheorem{remark}[theorem]{Remark}

\title{The Generalized Local Rule and the Burge--RSK Involution}
\author{Robin Langer}
\date{\today}

\begin{document}
\maketitle

\begin{abstract}
We analyse whether the Sch\"utzenberger involution relating the
Burge and RSK local rules extends to the $Q$-enriched setting
of~\cite{ADJOINT}.  The conclusion is negative: the $Q$-kernel
$\Phi_{a,b}(t)$ is context-sensitive (it depends on the strip
sizes at each face), and the Sch\"utzenberger involution changes
the interior partition shapes, propagating $Q$-weight changes
through the growth diagram.  The failure is structurally
informative: it shows that effective involutions for $Q$-positivity
must be \emph{local} (operating one face at a time), matching the
recursive structure of the shifted factorisation
$\Phi_{a,b}(t) = Q_a Q_b + t\,\Phi_{a+1,b+1}(t)$.
\end{abstract}

\tableofcontents


%% ============================================================
\section{Setup}
\label{sec:setup}
%% ============================================================

Recall the two parallel frameworks.

\medskip
\noindent\textbf{Algebraic.}
The fundamental commutation relation (cf.~\cite{COMMUTATION},
\cite{ADJOINT}):
\begin{equation}\label{eq:fund-comm}
  T(v)\, P(w) = \frac{1}{1-vw}\; P(w)\, T(v),
\end{equation}
which at the mode level gives the cascade relation:
\begin{equation}\label{eq:cascade}
  T_n\, P_m = \sum_{k=0}^{\min(n,m)} P_{m-k}\, T_{n-k}.
\end{equation}

\medskip
\noindent\textbf{Bijective.}
The Burge local rule $\mathfrak{U}_{\alpha,\beta}$ has type
signature
\begin{equation}\label{eq:type-sig}
  \mathfrak{U}_{\alpha,\beta}\colon
  \Z_{\geq 0} \times (D(\alpha) \cap D(\beta))
  \to U(\alpha) \cap U(\beta),
\end{equation}
where $U(\mu)$ (resp.\ $D(\lambda)$) is the set of partitions
obtained by adding (resp.\ removing) a horizontal strip.
The weight condition is:
\begin{equation}\label{eq:weight-cond}
  |\lambda| = |\alpha| + |\beta| - |\nu| + k.
\end{equation}

Each commutation step corresponds to one local rule application.
The scalar factor $1/(1-vw)$ generates the face labels.  Both
the Burge and RSK local rules share the same type
signature~\eqref{eq:type-sig} and satisfy the same weight
condition~\eqref{eq:weight-cond}.

\medskip
\noindent\textbf{The $Q$-enriched setting.}
In the $Q$-enriched framework (\cite{ADJOINT}, Part~III), the
face label $k$ carries weight $Q_{a+k}(q) \cdot Q_{b+k}(q)$,
where $a = |\alpha| - |\nu|$ and $b = |\beta| - |\nu|$ are the
horizontal strip sizes entering the diamond.  The face-label
generating function becomes the $Q$-kernel:
\begin{equation}\label{eq:Phi}
  \Phi_{a,b}(t) = \sum_{k \geq 0} Q_{a+k}\, Q_{b+k}\, t^k,
\end{equation}
replacing the universal $1/(1-t)$ of the standard rule.


%% ============================================================
\section{Defining generalised RSK}
\label{sec:gen-RSK}
%% ============================================================

This works without difficulty.  The RSK local rule has the same
type signature as Burge:
\begin{equation}\label{eq:RSK-type}
  \mathfrak{U}^{\mathrm{RSK}}_{\alpha,\beta}\colon
  \Z_{\geq 0} \times (D(\alpha) \cap D(\beta))
  \to U(\alpha) \cap U(\beta),
\end{equation}
with the same weight condition
$|\lambda| = |\alpha| + |\beta| - |\nu| + k$.

\begin{definition}[$Q$-enriched RSK]
\label{def:Q-RSK}
The \emph{$Q$-enriched RSK local rule} is defined identically to
the $Q$-enriched Burge rule: same partition geometry, same cascade
parameter~$k$, face label $k$ carries weight
$Q_{a+k}(q) \cdot Q_{b+k}(q)$.
\end{definition}

Since the inputs $(\nu, \alpha, \beta, k)$ are the same for both
rules, and the strip sizes $a$, $b$ depend only on the input, the
$Q$-weight at each individual face is the same for both rules
\emph{at that face considered in isolation}.


%% ============================================================
\section{The Sch\"utzenberger involution}
\label{sec:schutzenberger}
%% ============================================================

The composition
\begin{equation}\label{eq:psi-comp}
  \psi_R \circ \psi_B^{-1} \colon \CPP(\pi) \to \CPP(\pi)
\end{equation}
is well-defined: take a cylindric plane partition, decompose it
via the Burge bijection into a pair $(\gamma, \mathfrak{d})$
with $\gamma \in \IP$ and $\mathfrak{d} \in \ALCD(\pi)$, then
reassemble via the RSK bijection.  This preserves the base
partition~$\gamma$ and the face labels (the ALCD is unchanged).
In the standard (non-$Q$-enriched) case, this is a
weight-preserving endomorphism on $\CPP(\pi)$.

\begin{remark}
From~\cite{COMMUTATION}: ``Two natural local rules exist (RSK
and Burge), related by the Sch\"utzenberger involution.  Their
composition $\psi_R \circ \psi_B^{-1} \colon \CPP(\pi) \to
\CPP(\pi)$ has a conjectural interpretation via crystal bases
for $\widehat{\mathfrak{sl}}_n$.''
\end{remark}


%% ============================================================
\section{The propagation obstruction}
\label{sec:obstruction}
%% ============================================================

Here is the critical failure.  Consider a cylindric growth
diagram with multiple rows of faces.

\subsection{First row}

The inputs $(\nu, \alpha, \beta)$ at each first-row face are
determined by the base partition~$\gamma$.  These are the same
for both Burge and RSK.  Therefore:
\begin{itemize}
  \item Strip sizes at first-row faces: \textbf{identical} for
    both rules.
  \item $Q$-weights at first-row faces: \textbf{identical}.
  \item Output partition $\lambda$: \textbf{different}
    ($\lambda_B \neq \lambda_R$ in general).
\end{itemize}

\subsection{Second row and beyond}

At the second row, the inputs include~$\lambda$ from the first
row.  Since $\lambda_B \neq \lambda_R$, the strip sizes
$a'$,~$b'$ at second-row faces are \textbf{different} between
Burge and RSK.  Therefore:
\[
  Q_{a'+k} \cdot Q_{b'+k}
  \quad\text{differs between the two rules.}
\]

\begin{proposition}[Propagation obstruction]
\label{prop:obstruction}
The $Q$-weights propagate upward through the growth diagram.  The
Sch\"utzenberger involution changes the partition shapes at
interior vertices, which changes the strip sizes at subsequent
faces, which changes the $Q$-weights.  The total $Q$-weight
(the product over all faces) is \textbf{not} preserved.
\end{proposition}

Schematically:
\begin{center}
\begin{tabular}{l|l}
\textbf{Standard case} & \textbf{$Q$-enriched case} \\ \hline
Face weight $= 1$ (always) &
  Face weight $= Q_{a+k} \cdot Q_{b+k}$ \\[3pt]
Independent of strip sizes &
  Depends on strip sizes \\[3pt]
Sch\"utzenberger preserves &
  Sch\"utzenberger changes \\
total weight trivially &
  total $Q$-weight via propagation
\end{tabular}
\end{center}


%% ============================================================
\section{The $d = 1$ exception}
\label{sec:d1}
%% ============================================================

When $d = 1$, $Q_n = 1$ for all~$n$.  The $Q$-weight
$Q_{a+k} \cdot Q_{b+k} = 1$ regardless of $a$, $b$, $k$.  The
strip-size dependence vanishes, and the Sch\"utzenberger
involution trivially preserves $Q$-weights.  This recovers the
standard Burge--RSK equivalence.


%% ============================================================
\section{Context-sensitivity and global coupling}
\label{sec:context}
%% ============================================================

The failure is fundamentally because the $Q$-kernel is
\emph{context-sensitive}: $\Phi_{a,b}(t)$ depends on the strip
sizes $(a,b)$, which are determined by the global partition
structure.  In the standard case, the face-label generating
function is $1/(1-t)$ everywhere (universal), so any
rearrangement of the growth diagram preserves weights.  The
$Q$-enrichment breaks this universality.

\begin{remark}[Why positivity is hard]
\label{rem:hard}
This explains, at a structural level, why the positivity proof for
$Q_{n,c}(q)$ is hard: any bijective technique that rearranges the
growth diagram \emph{globally} (Sch\"utzenberger, evacuation, jeu
de taquin) will change the $Q$-weights unless it preserves the
exact strip sizes at every face.  The $Q$-kernel creates a
\textbf{global coupling} between the local rule and the partition
geometry.
\end{remark}


%% ============================================================
\section{What could work instead}
\label{sec:alternatives}
%% ============================================================

The failure of the global Burge--RSK approach suggests three
alternative strategies, all based on \emph{local} operations.


\subsection{Peel/glue at the ALCD level}

The sign-reversing involution from~\cite{FIRST_ATTEMPT} operates
on pairs $(S, \mathfrak{d})$ where $\mathfrak{d}$ is an ALCD.
The peel/glue mechanism uses Burge forward/inverse on a
\emph{single} face (the outermost diagonal), which does
\emph{not} change the rest of the growth diagram.  This is local
and self-consistent.  The $Q$-weight change at the peeled face is
controlled:
\begin{align*}
  \text{peel:} &\quad \text{removes one face with weight }
    Q_{a+k} \cdot Q_{b+k}, \\
  \text{glue:} &\quad \text{adds one face with weight }
    Q_{a+k} \cdot Q_{b+k}.
\end{align*}
The rest of the growth diagram is untouched, so its $Q$-weights
are preserved.  The involution pairs:
\[
  (S,\, \mathfrak{d})
  \;\longleftrightarrow\;
  (S \mathbin{\triangle} \{k\},\, \mathfrak{d}'),
\]
where $\mathfrak{d}'$ differs from $\mathfrak{d}$ by one face on
the outermost diagonal.


\subsection{The $Q$-kernel recurrence}

The shifted factorisation (\cite{ADJOINT}, Part~III,
Proposition~16.4):
\begin{equation}\label{eq:shifted}
  \Phi_{a,b}(t) = Q_a\, Q_b + t\, \Phi_{a+1,b+1}(t)
\end{equation}
peels off one layer of $Q$-weight at a time.  An involution that
peels/glues one diagonal at a time tracks the $Q$-weight shift
through this recurrence.  The $k = 0$ term $Q_a Q_b$ is the
``uncoupled'' contribution, and the remainder
$t\, \Phi_{a+1,b+1}(t)$ represents the cross-level interactions
that the involution must cancel.

Iterating~\eqref{eq:shifted}:
\begin{equation}\label{eq:shifted-iter}
  \Phi_{a,b}(t) = \sum_{j=0}^{N-1} Q_{a+j}\, Q_{b+j}\, t^j
    + t^N\, \Phi_{a+N,b+N}(t).
\end{equation}
So the involution can work layer by layer, cancelling one
cross-level term at a time.  This matches the Franklin-style
structure where each involution step moves exactly one element
between $S$ and the cylindric object.


\subsection{The commutator approach}

From~\cite{ADJOINT}, Part~III, the commutator isolates the
$k \geq 1$ terms:
\begin{equation}\label{eq:comm}
  [\mathbf{Q}(z),\, \mathbf{Q}^*(w)]
  = -\sum_{a,b \geq 0} P_a\, T_b\; z^a w^b\;
    \bigl(\Phi_{a,b}(zw) - Q_a\, Q_b\bigr).
\end{equation}
The obstruction to commutativity is precisely:
\begin{equation}\label{eq:obstruction}
  \Phi_{a,b}(t) - Q_a\, Q_b
  = t\, \Phi_{a+1,b+1}(t)
  = \sum_{k \geq 1} Q_{a+k}\, Q_{b+k}\, t^k.
\end{equation}
A successful involution must cancel precisely these cross-level
terms.  The peel/glue mechanism applied to the outermost face of
the ALCD amounts to cancelling the $k = 1$ term of this sum,
then iterating inward.


%% ============================================================
\section{The structural lesson}
\label{sec:lesson}
%% ============================================================

The Burge $\leftrightarrow$ RSK involution is a \emph{global}
transformation: it changes all internal vertices of the growth
diagram simultaneously.  But the $Q$-enrichment makes weights
\emph{locally coupled} to the partition shapes at those vertices.

\begin{theorem}[Locality principle]
\label{thm:locality}
Effective involutions for $Q$-positivity must be
\textbf{local}---operating on one diagonal/face at a time---because
the $Q$-kernel's shifted factorisation
\[
  \Phi_{a,b}(t) = Q_a\, Q_b + t\, \Phi_{a+1,b+1}(t)
\]
is itself a one-layer-at-a-time recurrence.
\end{theorem}

This is the algebraic reason why the Franklin-style peel/glue
approach (\cite{FIRST_ATTEMPT}) is the right framework: it matches
the recursive structure of the $Q$-kernel.  The involution peels
one face at a time, each step corresponding to one application of
the shifted factorisation, and the $Q$-weights are tracked through
the recurrence rather than through a global diagram transformation.


%% ============================================================
\section{Open question}
\label{sec:open}
%% ============================================================

Can the local peel/glue involution be made explicit?

\begin{conjecture}[Canonical peel]
\label{conj:peel}
Given $(S, \mathfrak{d})$ with $\mathfrak{d} \in
\ALCD_{n-j}(c)$, there exists a canonical choice of face on the
outermost diagonal to peel, satisfying:
\begin{enumerate}
  \item \emph{Involutivity.}\;
    $\mathrm{peel}(\mathrm{glue}(S, \mathfrak{d}))
    = (S, \mathfrak{d})$ and vice versa.
  \item \emph{Weight preservation.}\;
    $\sigma(S') + |\mathfrak{d}'|
    = \sigma(S) + |\mathfrak{d}|$.
  \item \emph{Fixed points.}\;
    The unpaired elements at $j = 0$ contribute exactly
    $Q_{n,c}(q)$.
\end{enumerate}
\end{conjecture}

The cyclic descent injection (\cite{NOTES}, Section~8) provides a
partial answer for the max-$1$ CPP case ($d \not\equiv 0
\pmod{3}$), but the general case---where the ALCD has arbitrary
depth---remains open.  The $Q$-kernel framework suggests that the
canonical choice should be guided by the shifted factorisation:
peel the face whose removal corresponds to the $k = 1 \to k = 0$
reduction in $\Phi_{a,b}(t)$.


%% ============================================================
\section{The pair-based $Q$-local rule}
\label{sec:pair-rule}
%% ============================================================

The preceding sections showed that na\"ively lifting the
Sch\"utzenberger involution to the $Q$-enriched setting fails
because of the propagation obstruction: partition diamonds couple
$Q$-weights to vertex shapes.  We now take a fundamentally
different approach.  The underlying combinatorial object for the
$Q$-commutator is \textbf{not} a partition diamond, but a
\textbf{pair of ground-set elements}.

\subsection{Why not diamonds}

The standard Burge/RSK local rule operates on a diamond
$(\nu;\, \alpha, \beta) \mapsto \lambda$ of integer partitions.
This is the bijective incarnation of $T_n P_m
= \sum_k P_{m-k} T_{n-k}$, which acts on partitions (the
Schur basis of~$\La$).

The $Q$-operator $\mathbf{Q}(z) = \sum Q_m z^m P_m$ packages
the cylindric Kostka polynomials as \emph{coefficients} of the
Pieri operators.  Each $Q_m$ is itself the signed generating
function of the ground set $\mathcal{G}_{m,c}$:
\begin{equation}\label{eq:Q-ground-set}
  Q_m = \sum_{(S,\mathfrak{d}) \in \mathcal{G}_{m,c}}
    (-1)^{|S|}\, q^{\sigma(S) + |\mathfrak{d}|}.
\end{equation}
The cascade $T_n P_m = \sum_k P_{m-k} T_{n-k}$ still happens
\emph{inside}, but the $Q$-weights are \emph{external} to the
cascade---they come from the ground-set decomposition.  A local
rule on diamonds captures the cascade but misses the cylindric
structure that $Q_m$ encodes.


\subsection{The underlying object: pairs of ground-set elements}

The product $\mathbf{Q}^*(w)\,\mathbf{Q}(z)$ involves terms
$Q_m\, Q_n\, z^m w^n\, T_n P_m$ (cf.~\cite{ADJOINT}, Part~II).
At the combinatorial level, the factor $Q_m \cdot Q_n$ is the
signed generating function of the \textbf{product ground set}:
\begin{equation}\label{eq:product-ground-set}
  Q_m \cdot Q_n
  = \sum_{\substack{(S,\mathfrak{d}) \in \mathcal{G}_{m,c} \\
    (S',\mathfrak{d}') \in \mathcal{G}_{n,c}}}
    (-1)^{|S|+|S'|}\,
    q^{\sigma(S) + |\mathfrak{d}| + \sigma(S') + |\mathfrak{d}'|}.
\end{equation}

After the cascade with $m = a + k$ and $n = b + k$, the coefficient
of $P_a\, T_b\, z^a w^b\, (zw)^k$ in
$\mathbf{Q}^*(w)\,\mathbf{Q}(z)$ is $Q_{a+k}\, Q_{b+k}$.

\begin{definition}[Pair-based $Q$-local rule]
\label{def:pair-rule}
The underlying combinatorial object at cascade level~$k$ with
strip sizes $(a,b)$ is a \emph{pair of ground-set elements}:
\begin{equation}\label{eq:pair}
  \bigl((S,\mathfrak{d}),\; (S',\mathfrak{d}')\bigr)
  \;\in\;
  \mathcal{G}_{a+k,c} \times \mathcal{G}_{b+k,c},
\end{equation}
with:
\begin{itemize}
  \item \textbf{Sign:}\; $(-1)^{|S|+|S'|}$.
  \item \textbf{$q$-weight:}\;
    $q^{\sigma(S) + |\mathfrak{d}|
      + \sigma(S') + |\mathfrak{d}'|}$.
  \item \textbf{Cascade parameter:}\; $k \geq 0$.
\end{itemize}
\end{definition}

The $Q$-kernel is then expressed as a signed sum over pairs:
\begin{equation}\label{eq:Phi-pairs}
  \Phi_{a,b}(t) = \sum_{k \geq 0} t^k
    \sum_{\substack{(S,\mathfrak{d}) \in \mathcal{G}_{a+k,c} \\
      (S',\mathfrak{d}') \in \mathcal{G}_{b+k,c}}}
    (-1)^{|S|+|S'|}\,
    q^{\sigma(S) + |\mathfrak{d}|
      + \sigma(S') + |\mathfrak{d}'|}.
\end{equation}


\subsection{The sign structure}

The sign $(-1)^{|S|+|S'|}$ depends on the \emph{total subset
size} $|S| + |S'|$, which ranges from $0$ to $(a+k) + (b+k)$.
The terms decompose as:

\begin{center}
\begin{tabular}{l|l|l}
$|S| + |S'|$ & \textbf{Sign} & \textbf{Contribution} \\ \hline
$0$ (both empty) & $+1$ &
  $\sum_{\mathfrak{d},\mathfrak{d}'}
  q^{|\mathfrak{d}| + |\mathfrak{d}'|}$ \\[3pt]
$1$ (one singleton) & $-1$ &
  cancels against $|S|+|S'| = 0$ terms \\[3pt]
$2$ & $+1$ &
  cancels against $|S|+|S'| = 1$ terms \\[3pt]
$\vdots$ & $\vdots$ & $\vdots$
\end{tabular}
\end{center}

\begin{remark}[Two sources of sign]
\label{rem:two-signs}
The signs arise from two \emph{independent} copies of the
alternating sum~\eqref{eq:Q-ground-set}: one from
$\mathcal{G}_{a+k}$ (the ``creation'' side) and one from
$\mathcal{G}_{b+k}$ (the ``annihilation'' side).  If each
factor $Q_{a+k}$ and $Q_{b+k}$ were independently non-negative,
the product would be non-negative.  The pair structure opens the
possibility of a \emph{coupled} involution that works even when
the individual factors are hard to handle.
\end{remark}


\subsection{The cascade as coupling}

The cascade parameter $k$ \textbf{couples} the two components of
the pair: the first component lives in $\mathcal{G}_{a+k,c}$ and
the second in $\mathcal{G}_{b+k,c}$, with $k$ shared.  This
coupling is absent in the standard Burge rule (where the face
label~$k$ carries weight~$1$ regardless of context) and is the
essential new feature of the pair-based framework.

An involution on pairs can exploit the shared cascade in ways that
two independent involutions cannot:
\begin{enumerate}
  \item \textbf{Cross-pair exchange.}
    Move an element from $S$ to $S'$ (or vice versa), adjusting
    the ALCD components to preserve total weight.  This changes
    the parity of $|S| + |S'|$ by $0$ (not useful for sign
    cancellation) but can redistribute weight between the two
    components.

  \item \textbf{Cascade shift.}
    Move from cascade level~$k$ to level~$k \pm 1$ by
    transferring data between the pair and the cascade.  This
    changes the ground-set levels from
    $(a+k, b+k)$ to $(a+k \pm 1, b+k \pm 1)$ and modifies the
    pair accordingly.

  \item \textbf{Coupled peel/glue.}
    In the first component, peel an element from
    $\mathfrak{d}$ and deposit it in~$S$ (changing the sign);
    simultaneously, in the second component, glue an element
    from $S'$ onto $\mathfrak{d}'$ (changing the sign back).
    The net sign change is $(-1)^{+1}\cdot(-1)^{-1} = 1$
    (no sign change), but the \emph{weight} can be redistributed.

  \item \textbf{Asymmetric peel.}
    Peel in only \emph{one} component: add an element to~$S$
    (sign changes by $-1$), adjusting the weight.  This is the
    analogue of the Franklin-style involution, but applied to one
    component of the pair.
\end{enumerate}

\begin{remark}[The cascade recurrence, revisited]
\label{rem:cascade-pair}
The shifted factorisation $\Phi_{a,b}(t) = Q_a Q_b + t\,
\Phi_{a+1,b+1}(t)$ acquires a new interpretation in the
pair framework.  The $k = 0$ term $Q_a Q_b$ corresponds to
pairs in $\mathcal{G}_{a,c} \times \mathcal{G}_{b,c}$ (no
cascade).  The remainder $t\, \Phi_{a+1,b+1}(t)$ corresponds
to pairs in $\mathcal{G}_{a+1,c} \times \mathcal{G}_{b+1,c}$
and higher (with cascade $k \geq 1$).  The involution must
map between these two regimes, using the cascade as the
transfer mechanism.
\end{remark}


%% ============================================================
\section{The involution programme on pairs}
\label{sec:pair-programme}
%% ============================================================

\subsection{The full signed sum}

Assembling all the pieces, the $Q$-kernel at cascade level~$k$ is
a signed sum over pairs:
\[
  [t^k]\, \Phi_{a,b}(t) = Q_{a+k}\, Q_{b+k}
  = \!\!\!\!
    \sum_{\substack{(S,\mathfrak{d}) \in \mathcal{G}_{a+k} \\
      (S',\mathfrak{d}') \in \mathcal{G}_{b+k}}}
    \!\!\!\!
    (-1)^{|S|+|S'|}\,
    q^{\sigma(S) + |\mathfrak{d}|
      + \sigma(S') + |\mathfrak{d}'|}.
\]
We need a sign-reversing, weight-preserving involution
$\iota$ on the set
\begin{equation}\label{eq:pair-set}
  \mathcal{P}_{a,b,k}
  := \mathcal{G}_{a+k,c} \times \mathcal{G}_{b+k,c}
\end{equation}
whose fixed points have $|S| + |S'|$ even and contribute
$Q_{a+k}\, Q_{b+k} \in \Z_{\geq 0}[q]$.

\subsection{Factored vs.\ coupled involutions}

The simplest approach is a \emph{factored involution}:
apply a sign-reversing involution $\iota_1$ on
$\mathcal{G}_{a+k}$ independently of a sign-reversing
involution $\iota_2$ on $\mathcal{G}_{b+k}$, and set
$\iota = \iota_1 \times \iota_2$.

\begin{proposition}\label{prop:factored}
If sign-reversing involutions $\iota_1$ on
$\mathcal{G}_{a+k,c}$ and $\iota_2$ on $\mathcal{G}_{b+k,c}$
exist (proving $Q_{a+k} \geq 0$ and $Q_{b+k} \geq 0$
individually), then $\iota = \iota_1 \times \iota_2$ is a
sign-reversing involution on $\mathcal{P}_{a,b,k}$.
\end{proposition}

\begin{proof}
If $(S,\mathfrak{d})$ is paired with $(T,\mathfrak{e})$ by
$\iota_1$ (so $|S| \not\equiv |T| \pmod{2}$), then for any
$(S',\mathfrak{d}')$, the pairs
$((S,\mathfrak{d}),(S',\mathfrak{d}'))$ and
$((T,\mathfrak{e}),(S',\mathfrak{d}'))$ have opposite signs and
the same weight.  The product involution pairs them.
\end{proof}

This reduces the pair problem to the original single-factor
problem and is therefore no easier.  The interest of the pair
framework lies in the possibility of a \textbf{coupled involution}
that uses the cascade to transfer data between components.

\subsection{Coupled involutions via cascade shift}

\begin{definition}[Cascade-coupled involution]
\label{def:cascade-coupled}
A \emph{cascade-coupled involution} on $\bigsqcup_{k \geq 0}
\mathcal{P}_{a,b,k}$ is a sign-reversing, weight-preserving
involution $\iota$ that may pair elements at \emph{different}
cascade levels:
\[
  \iota\colon
  \bigl((S,\mathfrak{d}),\, (S',\mathfrak{d}'),\, k\bigr)
  \;\longmapsto\;
  \bigl((T,\mathfrak{e}),\, (T',\mathfrak{e}'),\, k'\bigr),
\]
where $k' \neq k$ is allowed, provided that:
\begin{enumerate}
  \item $|S|+|S'| \not\equiv |T|+|T'| \pmod{2}$
    (sign reversal),
  \item $\sigma(S) + |\mathfrak{d}| + \sigma(S')
    + |\mathfrak{d}'| + k
    = \sigma(T) + |\mathfrak{e}| + \sigma(T')
    + |\mathfrak{e}'| + k'$
    (weight preservation, where the cascade contributes
    weight~$k$ via the $t^k$ monomial),
  \item $\iota^2 = \mathrm{id}$ (involutivity).
\end{enumerate}
\end{definition}

The cascade shift is the key new feature: an element at cascade
level~$k$ with pair data
$((S,\mathfrak{d}), (S',\mathfrak{d}'))$ can be paired with
an element at level~$k'$ by transferring weight between the
cascade parameter and the pair components.

\begin{remark}[The shifted factorisation as cascade shift]
The shifted factorisation
$\Phi_{a,b}(t) = Q_a Q_b + t\,\Phi_{a+1,b+1}(t)$ says: the
$k = 0$ pairs in $\mathcal{G}_a \times \mathcal{G}_b$ and the
$k \geq 1$ pairs in
$\bigsqcup_{k \geq 1} \mathcal{G}_{a+k} \times
\mathcal{G}_{b+k}$ must combine to give non-negative
coefficients.  A cascade-coupled involution can pair an element
at level~$k = 0$ with an element at level~$k = 1$ (or higher),
using the cascade parameter to absorb the weight difference.
\end{remark}

\subsection{The programme}

\begin{enumerate}
  \item \textbf{Define the pair set.}
    $\mathcal{P}_{a,b} := \bigsqcup_{k \geq 0}
    \mathcal{G}_{a+k,c} \times \mathcal{G}_{b+k,c}$, with sign
    $(-1)^{|S|+|S'|}$ and weight
    $\sigma(S) + |\mathfrak{d}| + \sigma(S') + |\mathfrak{d}'|
    + k$.

  \item \textbf{Construct a cascade-coupled involution}
    $\iota$ on $\mathcal{P}_{a,b}$.  The involution may pair
    elements at different cascade levels, using the cascade
    parameter as a ``bridge'' between the two ground-set
    components.

  \item \textbf{Identify the fixed points.}
    The unpaired elements of $\iota$ contribute
    $\Phi_{a,b}(t) \in \Z_{\geq 0}[q][\![t]\!]$, which implies
    $Q_{a+k} Q_{b+k} \in \Z_{\geq 0}[q]$ for each~$k$.

  \item \textbf{Deduce single-factor positivity.}
    Setting $b = 0$ gives
    $\Phi_{a,0}(t) = C_a(t) = \sum_k Q_{a+k} Q_k\, t^k$.
    Since $Q_0 = 1$, the $k = 0$ term is $Q_a$, and
    non-negativity of $\Phi_{a,0}(t)$ implies $Q_a \geq 0$.
\end{enumerate}

\begin{remark}[Advantage over the factored approach]
\label{rem:advantage}
Step~(4) is the key advantage of the pair framework: proving
$\Phi_{a,0}(t) \geq 0$ for all~$a$ is \emph{a priori} a
different problem from proving $Q_n \geq 0$ for all~$n$, and
the pair structure with cascade coupling provides additional
tools (cross-component transfer, cascade shifts) that are not
available when working with a single ground set.  In particular,
the involution on $\mathcal{G}_{a+k} \times \mathcal{G}_{k}$
can exploit the asymmetry between the two components (one at
level $a + k$, one at level~$k$) to construct pairings that
have no analogue in the single-factor setting.
\end{remark}


%% ============================================================
\section{The Q-Burge and Q-RSK rules}
\label{sec:two-rules}
%% ============================================================

The standard Burge and RSK local rules differ in a single
respect: the cascade direction applied to overlapping columns
in conjugate coordinates.
\begin{center}
\begin{tabular}{l|l}
\textbf{Rule} & \textbf{Cascade direction} \\ \hline
Burge forward & UP:
  $\varepsilon(i) = \min\{j > i : j \notin A \cup B \cup C\}$
  \\
RSK forward & DOWN:
  $\delta(i) = \max\{j < i : j \notin A \cup B \cup C\}$
\end{tabular}
\end{center}
Both have the same type signature and weight condition; they
produce different outputs only when the overlap
$A \cap B \neq \emptyset$.  We now lift this dichotomy to the
pair framework of Section~\ref{sec:pair-rule}.


\subsection{Peel and glue on ground-set elements}

Let $(S, \mathfrak{d}) \in \mathcal{G}_{n,c}$ with
$S \subseteq \{1, \ldots, n\}$, $|S| = j$,
$\mathfrak{d} \in \ALCD_{n-j}(c)$.  The \emph{peel} and
\emph{glue} operations transfer data between the subset~$S$
and the ALCD~$\mathfrak{d}$, using the local rule as the
transfer mechanism.

\begin{definition}[Peel]
\label{def:peel}
Suppose $\mathfrak{d}$ has a non-trivial outermost face.  The
\emph{$R$-peel} (for $R \in \{\text{Burge}, \text{RSK}\}$)
applies the inverse local rule~$R^{-1}$ to extract the
outermost face, producing:
\[
  (S, \mathfrak{d})
  \;\xrightarrow{\;\text{$R$-peel}\;}
  (S \cup \{\ell\},\; \mathfrak{d}'),
\]
where $\ell \in \{1, \ldots, n\} \setminus S$ is the position
determined by the inverse rule
(Burge$^{-1}$ cascades \textbf{down}, RSK$^{-1}$ cascades
\textbf{up}), and $\mathfrak{d}' \in \ALCD_{n-j-1}(c)$ is the
reduced ALCD\@.  Weight is preserved:
$\sigma(S \cup \{\ell\}) + |\mathfrak{d}'|
= \sigma(S) + |\mathfrak{d}|$, and the sign changes:
$(-1)^{j+1} = -(-1)^j$.
\end{definition}

\begin{definition}[Glue]
\label{def:glue}
The inverse of peel.  Given $i \in S$, the \emph{$R$-glue}
of~$i$ applies the forward local rule~$R$ to insert a
horizontal strip of weight~$i$ into $\mathfrak{d}$:
\[
  (S, \mathfrak{d})
  \;\xrightarrow{\;\text{$R$-glue of $i$}\;}
  (S \setminus \{i\},\; \tilde{\mathfrak{d}}),
\]
with the same weight preservation and sign reversal.
\end{definition}

In general, $\ell_{\text{Burge}} \neq \ell_{\text{RSK}}$ and
$\mathfrak{d}'_{\text{Burge}} \neq \mathfrak{d}'_{\text{RSK}}$:
the Burge inverse cascades the extracted data \emph{downward}
(tending to produce a smaller~$\ell$), while the RSK inverse
cascades it \emph{upward} (tending to produce a larger~$\ell$).


\subsection{The Q-Burge rule on pairs}

\begin{definition}[Q-Burge rule]
\label{def:Q-Burge}
On the pair set $\mathcal{P}_{a,b,k} = \mathcal{G}_{a+k,c}
\times \mathcal{G}_{b+k,c}$, the \emph{Q-Burge rule}
$\mathfrak{B}$ applies a Burge-peel to the \textbf{first
component}:
\[
  \mathfrak{B}\colon
  \bigl((S, \mathfrak{d}),\; (S', \mathfrak{d}')\bigr)
  \;\longmapsto\;
  \bigl((S \cup \{\ell_B\},\, \mathfrak{d}_B),\;
    (S', \mathfrak{d}')\bigr),
\]
where $\ell_B$ and $\mathfrak{d}_B$ are determined by
Burge$^{-1}$ applied to the outermost face of~$\mathfrak{d}$.
The second component is unchanged.
\end{definition}

\subsection{The Q-RSK rule on pairs}

\begin{definition}[Q-RSK rule]
\label{def:Q-RSK}
On the same pair set, the \emph{Q-RSK rule} $\mathfrak{R}$
applies an RSK-peel to the \textbf{second component}:
\[
  \mathfrak{R}\colon
  \bigl((S, \mathfrak{d}),\; (S', \mathfrak{d}')\bigr)
  \;\longmapsto\;
  \bigl((S, \mathfrak{d}),\;
    (S' \cup \{\ell_R\},\, \mathfrak{d}'_R)\bigr),
\]
where $\ell_R$ and $\mathfrak{d}'_R$ are determined by
RSK$^{-1}$ applied to the outermost face
of~$\mathfrak{d}'$.  The first component is unchanged.
\end{definition}

Both $\mathfrak{B}$ and $\mathfrak{R}$ are sign-reversing: each
changes $|S| + |S'|$ by exactly~$1$, so the total sign
$(-1)^{|S|+|S'|}$ flips.

\begin{remark}[The double asymmetry]
\label{rem:double-asymmetry}
The Q-Burge rule uses the Burge direction (cascade
\textbf{down} for the inverse) on the \emph{first} component,
while the Q-RSK rule uses the RSK direction (cascade
\textbf{up} for the inverse) on the \emph{second} component.
This double asymmetry---different cascade directions on
different components of the pair---is the structural feature
that makes the combined involution possible.  It mirrors the
standard Burge--RSK dichotomy, but applied to the two sides of
the pair rather than to two bijections on the same object.
\end{remark}


%% ============================================================
\section{The paired involution}
\label{sec:paired-involution}
%% ============================================================

Having two sign-reversing operations on the pair set---one for
each component, using opposite cascade directions---we combine
them into a single sign-reversing involution.


\subsection{The Franklin-style comparison}

In the classical Franklin involution (for the pentagonal number
theorem), one compares a \emph{peel candidate}~$\ell$ (the
smallest part of the partition) with a \emph{glue
candidate}~$g$ (the end of the staircase): if $\ell < g$, one
peels; if $g < \ell$, one glues; if $\ell = g$, the element
is a fixed point.

In the pair framework, the two rules provide two peel
candidates---one for each component---and the comparison
between them determines which component to act on.

\begin{definition}[The Q-Burge--RSK involution]
\label{def:Q-BR-involution}
Let $p = ((S, \mathfrak{d}),\, (S', \mathfrak{d}'))$ be a
pair in $\mathcal{P}_{a,b,k}$.  Define the involution
$\iota(p)$ as follows.

\medskip
\noindent\textbf{Step~1: Compute candidates.}
\begin{itemize}
  \item If $\mathfrak{d}$ is non-trivial, compute the
    Burge-peel candidate~$\ell_B$ for the first component
    (cascade \textbf{down}: $\ell_B$ is the position where the
    extracted face data lands).
  \item If $\mathfrak{d}'$ is non-trivial, compute the
    RSK-peel candidate~$\ell_R$ for the second component
    (cascade \textbf{up}: $\ell_R$ is the position where the
    extracted face data lands).
\end{itemize}

\medskip
\noindent\textbf{Step~2: Compare and act.}
\begin{enumerate}
  \item[(P)] \textbf{Peel phase.}\; If both candidates exist,
    compare $\ell_B$ and~$\ell_R$.
    \begin{itemize}
      \item If $\ell_B < \ell_R$: apply $\mathfrak{B}$ (peel
        first component).
      \item If $\ell_R < \ell_B$: apply $\mathfrak{R}$ (peel
        second component).
      \item If $\ell_B = \ell_R$: proceed to the glue phase.
    \end{itemize}
    If only one candidate exists, apply the corresponding rule.

  \item[(G)] \textbf{Glue phase.}\; If the previous step
    produced a peeled pair~$p'$, this phase does not apply.
    Otherwise, check whether a glue operation can reverse a
    previous peel:
    \begin{itemize}
      \item If $S \neq \emptyset$ and $\max(S)$ is
        Burge-glueable (i.e., Burge-glue of $\max(S)$
        followed by Burge-peel recovers~$\max(S)$): apply
        Burge-glue to the first component.
      \item Else if $S' \neq \emptyset$ and $\min(S')$ is
        RSK-glueable: apply RSK-glue to the second component.
    \end{itemize}

  \item[(F)] \textbf{Fixed point.}\; If neither peel nor glue
    applies, $p$ is a \textbf{fixed point} of~$\iota$.
\end{enumerate}
\end{definition}

The peel--glue duality ensures involutivity: if $\iota$ peels
the first component of~$p$ to produce~$p'$, then $\iota$
applied to~$p'$ glues the first component back (since the
newly added element $\ell_B \in S$ is the Burge-glue candidate,
and gluing it reverses the peel).


\subsection{Sign-reversal and weight preservation}

\begin{proposition}\label{prop:sign-weight}
The involution $\iota$ of Definition~\ref{def:Q-BR-involution}
is:
\begin{enumerate}
  \item \textbf{Sign-reversing:} each non-fixed-point pair~$p$
    satisfies
    $(-1)^{|S(\iota(p))| + |S'(\iota(p))|}
    = -(-1)^{|S(p)| + |S'(p)|}$,
    since exactly one of $|S|$, $|S'|$ changes by~$1$.
  \item \textbf{Weight-preserving:} the total $q$-weight
    $\sigma(S) + |\mathfrak{d}| + \sigma(S') + |\mathfrak{d}'|$
    is unchanged, since peel and glue individually preserve
    weight.
\end{enumerate}
\end{proposition}


\subsection{Fixed-point structure}

A pair $p = ((S, \mathfrak{d}), (S', \mathfrak{d}'))$ is a
fixed point of~$\iota$ when:
\begin{enumerate}
  \item Both ALCDs $\mathfrak{d}$ and $\mathfrak{d}'$ are
    trivial (no face to peel), \textbf{and}
  \item Both subsets $S$ and $S'$ are empty (no element to
    glue);
\end{enumerate}
or when the peel candidates coincide ($\ell_B = \ell_R$) and
no glue operation is available.

\begin{remark}[The trivial fixed points]
\label{rem:trivial-fp}
In the first case, $|S| = |S'| = 0$, so the fixed point has
sign $(-1)^0 = +1$ and contributes
$q^{|\mathfrak{d}| + |\mathfrak{d}'|}$ with a \emph{positive}
coefficient.  These trivial fixed points correspond to pure
ALCD pairs
$(\mathfrak{d}, \mathfrak{d}') \in \ALCD_{a+k}(c) \times
\ALCD_{b+k}(c)$ and contribute the ``uncancelled positive
core'' of $Q_{a+k} \cdot Q_{b+k}$.
\end{remark}


\subsection{Why two rules are needed}

\begin{remark}[Single rule vs.\ paired rules]
\label{rem:single-vs-paired}
A single-rule involution (say, Burge-peel/glue on the first
component alone) is a \emph{factored involution}
(Proposition~\ref{prop:factored}): it proves $Q_{a+k} \geq 0$
for the first component independently, which is the original
unsolved problem.

The paired involution uses \emph{both} rules simultaneously,
one per component, with opposite cascade directions.  The
cross-component comparison in Step~2(P) creates pairings that
have no analogue in the single-factor setting:
\begin{itemize}
  \item A pair where the first component is a fixed point of
    Burge-peel may still be paired by the Q-RSK rule acting on
    the second component.
  \item The opposite cascade directions (Burge \textbf{down}
    on first, RSK \textbf{up} on second) ensure that the peel
    candidates $\ell_B$ and $\ell_R$ explore different parts
    of the position space, maximising the chance that at least
    one operation is available.
\end{itemize}
\end{remark}


\subsection{Connection to the shifted factorisation}

The shifted factorisation
$\Phi_{a,b}(t) = Q_a Q_b + t\,\Phi_{a+1,b+1}(t)$ acquires a
concrete meaning in this framework.  The $k = 0$ pairs
$\mathcal{G}_a \times \mathcal{G}_b$ are acted on by the
involution~$\iota$ restricted to cascade level~$0$.  The
unpaired (fixed-point) pairs at level~$0$ contribute $Q_a Q_b$
to $\Phi_{a,b}(0)$.  At level~$k = 1$, the involution acts on
$\mathcal{G}_{a+1} \times \mathcal{G}_{b+1}$, and the
fixed-point contribution is $Q_{a+1} Q_{b+1}$, the coefficient
of~$t$.  In general:
\[
  [t^k]\, \Phi_{a,b}(t)
  = \sum_{\substack{p \in \mathcal{P}_{a,b,k} \\
    \iota(p) = p}}
  q^{\mathrm{wt}(p)}
  \;\in\; \Z_{\geq 0}[q],
\]
provided $\iota$ has only positive fixed points at each cascade
level.  Setting $b = 0$ recovers $Q_a$ as the $k = 0$
fixed-point contribution of the involution on
$\mathcal{G}_{a} \times \mathcal{G}_{0}$, where
$\mathcal{G}_0 = \{(\emptyset, \mathfrak{d}_0)\}$ is a
singleton (the trivial ALCD at level~$0$).


\begin{thebibliography}{99}
\bibitem{ADJOINT}
  R.~Langer,
  \emph{The $Q$-Operator and its Adjoint: From Level~3 Operator
  Traces to Cascade Commutators},
  companion note, 2026.

\bibitem{Borodin}
  A.~Borodin,
  \emph{Periodic Schur process and cylindric partitions},
  Duke Math.~J.\ \textbf{140} (2007), 391--468.

\bibitem{COMMUTATION}
  R.~Langer,
  \emph{From Commutation Relations to Growth Diagrams:
  A Summary of the Operator--Bijection Correspondence},
  companion note.

\bibitem{FIRST_ATTEMPT}
  R.~Langer,
  \emph{Towards an Involution Proof of Positivity for $Q_{n,c}(q)$},
  companion note.

\bibitem{NOTES}
  R.~Langer,
  \emph{Weighted Local Rules and the Cyclic Descent Injection},
  companion note.

\bibitem{UNDERLYING_SET}
  R.~Langer,
  \emph{The Underlying Set: Cylindric Plane Partitions, ALCDs,
  and the Ground Set for the Involution},
  companion note.

\bibitem{langer-thesis}
  R.~Langer,
  \emph{Cylindric Partitions and Bijective Proofs},
  PhD thesis, University of Melbourne.

\bibitem{vanLeeuwen}
  M.~A.~A.~van~Leeuwen,
  \emph{The Robinson--Schensted and Sch\"utzenberger algorithms,
  an elementary approach},
  Electron.~J.~Combin.\ \textbf{3} (1996), R15.
\end{thebibliography}

\end{document}
